\documentclass{article}
\usepackage{pxLST}
\begin{document}

\section*{pxLST float test}

This document exercises the \texttt{pxCodeBox} and \texttt{pxCodeBreak} environments
with captions and cross-references.

\subsection*{1. pxCodeBox with label for cross-reference}
\pxLSTsetup{style=dark, language=pxJava, numbers=left}
\begin{pxCodeBox}{float1}{Hello World}{[lst:float1]}{A simple Java Hello World listing.}
public class HelloWorld {
    public static void main(String[] args) {
        System.out.println("Hello, World!");
    }
}
\end{pxCodeBox}

See listing tagged \texttt{[lst:float1]} above.

\subsection*{2. pxCodeBox with Bash language}
\pxLSTsetup{style=framed, language=pxBash, numbers=none}
\begin{pxCodeBox}{float2}{Build Script}{[lst:float2]}{Bash build script for P3CTeX.}
#!/bin/bash
cd tex/
pdflatex -halt-on-error main.tex
echo "Build complete."
\end{pxCodeBox}

See the build script in listing \texttt{[lst:float2]}.

\subsection*{3. pxCodeBreak (breakable listing)}
\pxLSTsetup{style=dark, language=algoritme, numbers=left}
\begin{pxCodeBreak}{float3}{Factorial Algorithm}{[lst:float3]}{Iterative factorial in pseudocode.}
Funcio factorial(n: Integer): Integer
  Si n <= 0 Llavors
    Retornar 1
  Sinó
    result <-- 1
    Per i = 1 A n Fer
      result <-- result * i
    Fi Per
    Retornar result
  Fi Si
Fi Funcio
\end{pxCodeBreak}

See the algorithm listing tagged as \texttt{[lst:float3]}.

\subsection*{4. pxCode (simple, no caption)}
\pxLSTsetup{style=plain, language=pxJava, numbers=none}
\begin{pxCode}{simple-inline}
// inline, no caption
int x = 42;
\end{pxCode}

\subsection*{5. pxCodeBox with light style}
\pxLSTsetup{style=light, language=pxYAML, numbers=none}
\begin{pxCodeBox}{float4}{YAML Config}{[lst:float4]}{YAML configuration example.}
server:
  host: localhost
  port: 8080
  debug: false
\end{pxCodeBox}

\subsection*{6. Multiple listings cross-referencing each other}
Listing \texttt{[lst:float1]} shows Java.
Listing \texttt{[lst:float2]} shows Bash.
Listing \texttt{[lst:float3]} shows an algorithm.
Listing \texttt{[lst:float4]} shows YAML.

\end{document}
